\documentclass{article}
\usepackage[utf8]{inputenc}
\usepackage[T1]{fontenc}
\usepackage{geometry}
\usepackage{graphicx}
\usepackage{float}
\usepackage{hyperref}
\usepackage[english,indonesian]{babel}

\title{Optimasi Pemrosesan Template LaTeX Menggunakan Pendekatan Berbasis AI}
\author{Academic Support Agent}
\date{\today}

\begin{document}
\maketitle

\begin{abstract}

\end{abstract}

\chapter{PENDAHULUAN}


\chapter{TINJAUAN PUSTAKA}


% Original content (commented out):
% \documentclass{article}
\usepackage[utf8]{inputenc}
\usepackage{float}
\usepackage{graphicx}
\graphicspath{images}
\usepackage[style=apa]{biblatex}
\addbibresource{bib.bib}
\usepackage{geometry}
 \geometry{
 a4paper,
 total={170mm,257mm},
 left=20mm,
 top=20mm,
 }

% Replace (Title), (Name), and (Count) with relevant information
% Individual sections can be found in the Sections folder
% Insert images into the images folder
% Insert references to the bib.bib file


\title{(Title) \\[1ex] \large Project Proposal}
\author{(Name)}
\date{November 2022}

\begin{document}

\maketitle

\begin{center}
    School of Computer Science\\
    University of Lincoln\\
    United Kingdom

    \vspace{0.8cm}

    Word Count: (Count)
\end{center}

\tableofcontents
\pagebreak

\input{Sections/Section 1 - Introduction.tex}

\input{Sections/Section 2 - Aims and Objectives.tex}

\input{Sections/Section 3 - Project Plan and Risk Analysis.tex}


\section{References}
\printbibliography[heading=none]

\end{document}


\end{document}